\section{The proof chain closed so far}
\label{sec:proved-chain}

\begin{figure}[ht]
\centering
\begin{tikzpicture}[node distance=4mm]
  \node[chainnode] (kg) {Terminating internal KeyGen\\\(\prefix_{992}(sk)=vk\)};
  \node[chainnode,right=of kg] (mem) {Constructed Verify memory\\\(\operatorname{SKMemPrefix}\)};
  \node[chainnode,right=of mem] (mu) {Actual raw \(\mu\) procedures\\\(\take_{32}(\mu_S^{64})=\mu_V^{32}\)};
  \node[chainnode,right=of mu] (chal) {Actual \(\mu_{32}\) absorb\\Equal challenge suffix states};
  \draw[chainarrow] (kg) -- (mem);
  \draw[chainarrow] (mem) -- (mu);
  \draw[chainarrow] (mu) -- (chal);
  \node[provennode,below=12mm of kg] (apihelper) {Actual raw API\\export/trace/frame \statusproved};
  \node[provennode,right=7mm of apihelper] (apicomp) {Region-local \(\mu\) and\\exact sequence \statusproved};
  \node[partialnode,below=12mm of apicomp] (direct) {Stored \(\mu\) transport\\product/replay\\\statuspartial\\\statusdeferred};
  \node[opennode,below=12mm of chal] (full) {Highbits/LSB and full challenge\\\statusblocked};
  \draw[chainarrow,draw=provedcolor] (kg) -- (apihelper);
  \draw[chainarrow,draw=provedcolor] (apihelper) -- (apicomp);
  \draw[chainarrow,dashed,draw=partialcolor] (apicomp) -- (direct);
  \draw[chainarrow,dashed,draw=blockedcolor] (chal) -- (full);
\end{tikzpicture}
\caption{The verified raw/internal chain (above) and the Week~6/7 API semantic
boundary (below). This security boundary remains unchanged in Weeks~8--9.}
\label{fig:proof-chain}
\end{figure}

\subsection{Step 1: the packed-key prefix returned by KeyGen}

Jasmin \identifier{_pack_sk_m23} first copies \identifier{vkp[0..vkbytes)} to
\identifier{skp}. In mode~2, \identifier{vkbytes=992}, and subsequent writes
follow this layout.
\[
\underbrace{[0,992)}_{vk}
\;\Vert\;
\underbrace{[992,1184)}_{3\cdot64\text{-byte }\eta\text{ vector}}
\;\Vert\;
\underbrace{[1184,1376)}_{2\cdot96\text{-byte }\eta_2\text{ vector}}
\;\Vert\;
\underbrace{[1376,1408)}_{32\text{-byte key}}.
\]
\identifier{pack_vec_eta_to_prefix_frame} and
\identifier{pack_vec2_eta_to_prefix_frame} prove that later writes do not change
\([0,992)\).

\begin{verified}[Generated packer prefix]
\label{thm:packer-prefix}
\identifier{pack_sk_m23_mode2_vk_prefix} proves that, under mode-2 constants,
if the actual generated \identifier{_pack_sk_m23} terminates, then
\[
  \operatorname{VKPrefix}_{992}(res,vk_0)
\]
holds. Through \identifier{keypair_full_m23_mode2_return_prefix} and
\identifier{keypair_internal_mode2_return_prefix}, this result lifts to the
return values of actual generated \identifier{_keypair_full_m23} and
\identifier{crypto_sign_keypair_internal_mode2_jazz}.
\end{verified}

This theorem holds without analyzing the distributions inside every KeyGen retry
attempt: whichever attempt terminates, the final packer call produces the same
prefix. Precisely for this reason, it is a strong return-value invariant but
remains partial correctness.

\subsection{Step 2: from local arrays to Verify memory}

\identifier{Mode2KeyMemoryBridge.stored_vk_buffer_at} constructs concrete memory
by writing the verification key's first 992 bytes over existing memory using
\identifier{stores}. \identifier{keygen_prefix_memory_relation_is_constructible}
uses \cref{def:vk-prefix} to show that this memory satisfies
\cref{def:mem-prefix}. \identifier{keygen_prefix_reaches_mu_memory} converts it
to the exact predicate of the hash theorem.

\begin{verified}[Non-vacuity of raw-hash premises]
\label{thm:memory-witness}
\identifier{generated_raw_mu_preconditions_from_keygen_prefix} chooses
\(base=0\), \(prelen=0\), and stored verification-key memory to construct a
concrete witness satisfying the no-wrap, raw-branch, valid-region, and
\identifier{sk_memory_prefix} premises simultaneously.
\end{verified}

This result shows that the premises are consistent and reachable from internal
KeyGen return values. This witness is historical local evidence constructed with
an authored \identifier{stores} operation. The Week~5 API-level theorem reaches
the same relation using only actual caller memory and generated copy procedures,
without using this witness.

\subsection{Step 3: from copy helpers to the actual raw KeyGen caller}

\sourcepath{easycrypt/refinement/composition/ApiKeyMemoryBridge.ec} directly
targets actual generated procedures obtained by focused extraction: KeyGen's
\identifier{__kp_api_copy_2080_to_addr} and
\identifier{__kp_api_copy_2752_to_addr}, and Sign's and Verify's
\identifier{_api_copy_raw_to_2752_prefix}.

\begin{verified}[Prefix preservation by individual API copy helpers]
\label{thm:api-copy-prefix}
\identifier{keygen_export_vk_mode2_prefix} and
\identifier{keygen_export_sk_mode2_prefix} prove that terminating actual KeyGen
exporters write the 992-byte VK and 1408-byte SK, respectively, to raw memory.
\identifier{keygen_export_sk_mode2_prefix_frames_vk} additionally proves that
the actual SK exporter preserves the existing VK prefix under the minimal
\identifier{disjoint_regions} premise.
\identifier{sign_import_mode2_sk_prefix} and
\identifier{verify_import_mode2_vk_prefix} prove that the actual importers read
1408 bytes and 992 bytes, respectively, from valid input regions, and that the
first 992 bytes of both local arrays agree with the original memory.
\identifier{sign_verify_api_importer_same_inputs} also shows that the two
importers return equal values and memory on equal inputs.
\end{verified}

The pure adapter theorems \identifier{exported_mode2_regions_agree},
\identifier{imported_prefixes_agree}, and
\identifier{imported_sign_sk_reaches_mu_memory} provide logical bridges from
prefix agreement of the two exported regions to local Sign/Verify prefix
agreement and raw \(\mu\) memory premises.

In \sourcepath{RawApiAddressBridge.ec},
\identifier{canonical_ui64_address_roundtrip},
\identifier{mode2_vk_region_bridge}, and \identifier{mode2_sk_region_bridge}
prove that canonical integer addresses are unchanged through
\identifier{W64.of_int} and \identifier{W64.to_uint} and satisfy both
valid-region predicates. Thus the earlier exporter/importer address-representation
difference is closed under explicit canonical-region contracts.

In \sourcepath{RawApiKeygenExportComposition.ec},
\identifier{keypair_raw_api_exact_export_trace} first establishes equality of
final memory and return values between the actual raw KeyGen caller and a trace
making the two copies explicit.

\begin{verified}[Sequential external export by actual raw KeyGen]
\label{thm:raw-keygen-export}
\identifier{keypair_raw_api_exports_matching_prefixes} opens actual generated
\identifier{cryptolab_haetae_mode2_keypair_internal}. For every terminating
execution satisfying the premises of a 32-byte seed read, canonical regions for
the 992-byte VK and 1408-byte SK, and separation of the two outputs, it proves
that the return value is 0 and
\[
  \operatorname{ExternalKeyPrefix}(mem_{\mathrm{out}},sku,vku)
\]
holds. In other words, the VK exported first by the actual caller is preserved
by the subsequent SK export, and the first 992 bytes of the two external buffers
are pointwise equal.
\end{verified}

This advances composition of helper facts into one caller postcondition, but
remains partial correctness. It does not prove KeyGen retry termination, and
canonical/valid regions are calling contracts, not substitutes for physical
memory safety.

\subsection{Step 4: actual raw Sign/Verify \texorpdfstring{\(\mu\)}{mu} procedures}

\sourcepath{easycrypt/refinement/composition/ExtractedMuHashPrefix.ec} compares
the following generated helpers step by step.

\begin{itemize}
  \item Sign packed-SK absorb and Verify raw VK-memory absorb;
  \item The two raw absorbs for equal pre/message addresses and lengths;
  \item Keccak state initialization, permutation, and finalize;
  \item The first 32 bytes of Sign's 64-byte squeeze and Verify's 32-byte squeeze.
\end{itemize}

The resulting \identifier{sign_verify_raw_mu_top_wrappers_prefix} proves
\cref{def:mu-prefix} at the local wrapper level. The important Week~3 advance
was to remove wrappers from the conclusion surface.

\begin{verified}[Relation between actual generated raw \(\mu\) procedures]
\label{thm:generated-mu}
The two procedures of \identifier{sign_verify_generated_raw_mu_prefix} are
actual generated \identifier{Sign._sf_mu_rawpre} and
\identifier{Verify.__verify_hash_mu}. Under equal pre/message addresses and
lengths, \identifier{RawPreLen}, two \identifier{ValidRegion} predicates,
992-byte \identifier{SKMemPrefix}, and address non-wrapping, it proves
\[
  \operatorname{MuPrefix}_{32}(res_S,res_V)
\]
as the result relation.
\end{verified}

The Sign side is connected exactly to its wrapper by
\identifier{sf_mu_rawpre_refines_raw_top}, and the Verify side by
\identifier{verify_hash_mu_raw_refines_raw_top}; relational transitivity then
composes the intermediate theorem. Verify's raw branch is not assumed through
a separate flag. \identifier{raw_prelen_shr63_zero} proves that the actual
\identifier{prelen >> 63} computation yields 0, selecting that branch.

\subsection{Step 5: equality of the challenge suffix}

The Sign challenge transcript absorbs \(\take_{32}(\mu_S^{64})\) after the
64-byte highbits/LSB part, and Verify absorbs \(\mu_V^{32}\) at the same position.
The current theorem does not yet prove equality of those first 64 bytes;
it assumes equal challenge states and starting position 64.

\begin{verified}[Generated raw \(\mu\)-to-challenge composition]
\label{thm:challenge-suffix}
\identifier{generated_raw_mu_to_challenge_suffix_zero_loss} compares two
composed programs that execute actual \identifier{_sf_mu_rawpre} and actual
\identifier{__verify_hash_mu}, respectively, followed by the actual Sign/Verify
\(\mu_{32}\) absorb helpers. Under the premises of \cref{thm:generated-mu} and
initial challenge position 64, the returned challenge states and positions are
exactly equal.
\end{verified}

\identifier{TranscriptBytes.challenge_input_eq_components} supplies a pure
list lemma showing that equality of highbits and LSB bytes, once established,
yields equality of the entire input lists. The interface between the current
result and the next unresolved obligation is therefore explicit.

\subsection{Step 6: exact traces of actual raw Sign/Verify callers}

\sourcepath{RawApiCallerMuTrace.ec} exactly compares actual
\identifier{cryptolab_haetae_mode2_signature_internal} with a Sign trace
augmented by ghost observations. \identifier{sign_raw_api_exact_mu_trace}
preserves final memory and return values;
\identifier{sign_raw_api_trace_imports_mode2_sk} and
\identifier{sign_raw_api_trace_binds_hash_inputs} record the imported 992-byte
prefix and actual pre/message pointers and lengths.

On the Verify side, \sourcepath{RawApiVerifyMuTrace.ec} connects the full M23,
mode~2, internal, raw, and cryptolab wrappers exactly to their traces.
Crucially, Verify may reject before hashing.

\begin{verified}[Verify success reaches the actual hash tail]
\label{thm:verify-accept-tail}
\identifier{verify_raw_api_actual_accept_implies_trace_tail_reached} and
\identifier{verify_cryptolab_actual_accept_implies_trace_tail_reached} prove
that the trace's \identifier{tail_reached} is true when the actual return value
is 0. \identifier{verify_raw_api_actual_accept_binds_hash_inputs} and its
cryptolab counterpart bind the VK/pre/message pointers and lengths supplied to
actual \identifier{__verify_hash_mu} on that successful path to the caller arguments.
\end{verified}

This does not prove tail reachability for arbitrary signatures or acceptance of
legitimate Sign outputs. Rejecting executions need not reach hashing; tail
reachability of legitimate outputs is a separate functional-correctness obligation.

\subsection{Step 7: accepted-path generated adapters and the remaining gap}

In \sourcepath{RawApiMuReachability.ec},
\identifier{raw_api_region_reaches_mu_zero_loss} constructs
\cref{def:mem-prefix} from the actual external prefix, Sign import, and canonical
VK region. There is no artificial \identifier{stores mem 0} witness here.

\begin{verified}[Accepted traces instantiate generated hash/suffix theorems]
\label{thm:accepted-api-adapter}
\identifier{raw_api_accepting_execution_hash_mu_zero_loss} proves
\(\operatorname{MuPrefix}_{32}\) for the results of actual generated
\identifier{_sf_mu_rawpre} and \identifier{__verify_hash_mu}, under observations
supplied by the Sign trace and a successful Verify trace.
\identifier{raw_api_accepting_execution_generated_challenge_suffix_zero_loss}
also equates the actual generated suffix-helper results under the same premises
and challenge position 64.
\end{verified}

However, the program surfaces of these two theorems are the generated hash and
suffix helpers, not the trace modules themselves. Week~6 closed the Sign-output
frame and region-local memory relations below, but a separate semantic edge
transporting the relation to stored results of completed traces is still needed.

\subsection{Step 8: Sign-output frames and region-local sequential traces}

\sourcepath{RawApiSignOutputFrame.ec} opens the actual output copy and raw Sign
caller. \identifier{sign_output_copy_exact_and_frames_reused} writes the
1474-byte signature prefix exactly while framing bytes outside it, and
\identifier{sign_raw_api_frames_reused_regions} lifts this to actual
\identifier{cryptolab_haetae_mode2_signature_internal}.

\begin{verified}[Reused-region frame for actual Sign output]
\label{thm:sign-output-frame}
If the valid signature/siglen output regions are explicitly disjoint from the
VK, SK, pre, and message regions, then every terminating execution returns 0
from Sign, preserves the 992-byte VK, 1408-byte SK, and pre/message regions,
and writes 1474 to siglen.
\identifier{sign_raw_api_preserves_external_key_prefix} also transports
KeyGen's \identifier{ExternalKeyPrefix} across the actual Sign call. This
disjointness contract has a concrete witness, but signing-loop termination is
not a conclusion.
\end{verified}

\sourcepath{RegionLocalMuEquivalence.ec} and
\sourcepath{RegionLocalMuTop.ec} remove the whole-memory equality requirement.

\begin{verified}[Region-local equivalence of actual generated hashes]
\label{thm:region-local-mu}
\identifier{sign_verify_generated_raw_mu_prefix_regionwise} compares actual
\identifier{_sf_mu_rawpre} and \identifier{__verify_hash_mu}. It proves
\(\operatorname{MuPrefix}_{32}(res_S,res_V)\) using only the mode-2 SK/VK prefix
relation, raw branch and canonical/no-wrap regions, and
\cref{def:region-eq} for pre/message read regions. Unrelated memory bytes need
not agree.
\end{verified}

Finally, \identifier{raw_sign_then_verify_actual_exact_trace} in
\sourcepath{RawApiDirectObservedMu.ec} proves that a program executing actual
raw Sign followed by actual cryptolab Verify and the sequential execution of
the two traces preserve both results and final memory.
\identifier{sign_then_verify_trace_accept_implies_tail_reached} and
\identifier{sign_then_verify_accept_binds_observed_inputs} supply tail
reachability of successful Verify and both hash descriptors as postconditions.

Frames, byte locality, exact sequence, accept control, and input binding now all
compile. Nevertheless, the remaining \claim{OBL-DIRECT-OBSERVED-MU} cannot yet
directly compare \identifier{SignRawApiMuTrace.observed_mu} and
\identifier{VerifyCryptolabMuTrace.observed_mu}. This requires
\claim{OBL-MU-PRODUCT-REPLAY}, transporting the existing pRHL hash-return theorem
to ghost results of completed traces with different residual continuations.
Week~7 could not construct this without a new axiom or wrapper widening and
froze it as \statusdeferred. Thus
\(\delta_{\mu,\mathrm{raw\_api\_accept}}=0\) is still not recorded.

\subsection{Step 9: the actual mode-2 signature prefix codec}

The Week~7 functional-correctness path switched to the actual extracted
\identifier{_pack_sig_prefix} and \identifier{_unpack_sig_prefix}.
\identifier{pack_sig_prefix_mode2_layout} in
\sourcepath{Mode2SignaturePrefixPack.ec} proves that the 256 challenge bits
occupy the first 32 bytes and the low signed bytes of the 1024 low coefficients
occupy the next 1024 bytes. \identifier{unpack_sig_prefix_mode2_layout} in
\sourcepath{Mode2SignaturePrefixUnpack.ec} reads this layout in reverse and
preserves the unused low-array region \([1024,2048)\).
\sourcepath{Mode2SignaturePrefixRoundTrip.ec} defines an actual-procedure harness
calling the two procedures in sequence.

\begin{verified}[Signature prefix codec]
\label{thm:signature-prefix-codec}
The theorem \identifier{pack_unpack_sig_prefix_mode2_roundtrip} calls the two
actual generated procedures sequentially. Under
\cref{def:signature-prefix-canonical}, it proves that the 256 decoded challenge
entries and 1024 mode-2 low entries equal their inputs, and that the remaining
1024 words of decoded low are preserved. All-zero arrays show non-vacuity of
the premises.
\end{verified}

This is Hoare partial correctness for the 1056-byte prefix. It excludes size
metadata after byte 1056, hbz/h rANS payloads, canonical zero padding, the pack
output tail, and the challenge-array tail frame. It concludes neither the full
1474-byte codec nor \(\delta_{\mathrm{Encoding}}=0\).

\subsection{Step 10: actual HBZ leaves, table certificates, and pure rANS steps}

Instead of inlining \claim{OBL-SIG-HBZ-ENCODE-DECODE} all at once, Week~8
split it into four layers: coefficient/symbol leaves, tables, rANS core, and
full wrappers. \sourcepath{Mode2HbzCodecSpec.ec} fixes the mode-2 constants
\[
  \mathsf{count}=1024,\qquad m=13,\qquad \mathsf{offset}=6,\qquad
  \mathsf{scale}=2^{10},\qquad L=2^{23}
\]
and the predicates of \cref{def:canonical-hbz}.

\begin{verified}[Actual HBZ prepare/apply leaf inverse]
\label{thm:hbz-leaf-inverse}
\identifier{encode_hb_z1_prepare_mode2_correct} proves that actual
\identifier{encode_hb_z1_prepare_jazz} maps canonical coefficients
\(h_i\in[-6,7)\) to symbols \(h_i+6\), sets the bad word to 0, and preserves
the symbol tail. \identifier{decode_hb_z1_apply_mode2_correct} proves that the
actual apply procedure restores the original 32-bit coefficients from the same
symbols and preserves the coefficient tail.
\identifier{encode_prepare_decode_apply_mode2_inverse}, which executes the two
actual wrappers sequentially, concludes exact recovery of the first 1024
coefficients. The all-zero HBZ array shows non-vacuity of the canonical premises.
\end{verified}

This theorem is a local prepare/apply inverse; the rANS buffer, encoder size,
and full wrapper do not appear on its program surface. Separately,
\sourcepath{Mode2HbzActualBoundary.ec} establishes by exact equivalence that
the actual \identifier{_encode_hb_z1_full}/\identifier{_decode_hb_z1_full}
bodies in the focused Week~8 extraction and Week~7 signature pack/unpack
extractions are the same procedures. Thus there is no issue of the leaf theorems
and full packer referring to different code copies, but procedure identity
alone does not supply an inverse.

\begin{verified}[Actual table certificate and pure rANS inverse]
\label{thm:hbz-rans-step}
\leavevmode\par
\sourcepath{Mode2HbzTableCertificate.ec} shows that the 13 frequencies are
positive, sum to 1024, and have intervals partitioning \([0,1024)\). It proves
that the actual fast encoder expression computed using reciprocal
multiply/high-half/shift and bias equals mathematical \(E_s\) within the allowed
state range, without 32/64-bit overflow.
\identifier{actual_mode2_hbz_tables_certified} in
\sourcepath{Mode2HbzSymbolWordsGenerated.ec} applies
\cref{def:hbz-table-certificate} to the literal encoder/decoder/symbol arrays of
actual pack/unpack. Finally, \identifier{pure_rans_step_inverse} and
\identifier{hbz_fast_step_decode_inverse} in \sourcepath{Mode2RansCore.ec} prove
\[
  D_s(E_s(x))=x,\qquad
  0\le s<13,\quad 1\le x<\mathsf{hbz\_xmax}(s)
\]
with \(s=0,x=1\) as a concrete witness to the fast-step premises.
\end{verified}

\subsection{Step 11: rANS byte stacks, actual suffix copying, and harness control}

Instead of immediately connecting the encoder and decoder in lockstep, Week~9
first fixed the shared \(xs/cuts/bytes\) semantics of \cref{def:rans-byte-trace}.

\begin{verified}[Pure normalization traces and word arithmetic]
\label{thm:rans-byte-stack}
\identifier{renorm_len_le2} and \identifier{renorm_reduced_bounds} in
\sourcepath{Mode2RansByteStack.ec} show that, for \(2^{23}\le x<2^{31}\), each
symbol emits only 0, 1, or 2 bytes and the reduced state satisfies
\(1\le r<\mathsf{hbz\_xmax}(s)\).
\identifier{renorm_bytes_readback} proves that reading back the segment in
decoder order recovers the pre-normalization state;
\identifier{serialize32_parse_inverse} proves the round trip for the first
4-byte little-endian state.
\identifier{normalized_fast_step_state_bounds} and
\identifier{encode_trace_state_bounds} show that every boundary state of a
canonical trace remains in \([2^{23},2^{31})\), and
\identifier{valid_rans_trace_canonical} proves that trace objects are
constructed from their definitions. Separately,
\sourcepath{Mode2RansNormalization.ec} connects actual W32 shifts, low-byte
truncation, decoder shift/bitwise-or, and W64 cursor decrements to the same
integer expressions. In particular,
\identifier{encoder_two_byte_decoder_order} verifies that backward stores
produce the forward segment \([\mathsf{high},\mathsf{low}]\).
\end{verified}

These theorems are pure/word leaves that do not call the actual generated
encoder/decoder. They therefore do not justify inserting
\identifier{valid_rans_trace} as a precondition on actual output or assuming
decoder success.

\begin{verified}[Actual suffix copying and branching on encoder results]
\label{thm:rans-actual-control}
\identifier{copy_encoded_suffix_correct} proves for actual
\identifier{__copy_encoded_suffix} that, if
\[
  0\le off_0,\quad 0\le n,\quad off_0+n\le2048
\]
then \(out[j]=enc[off_0+j]\) for \(0\le j<n\), and the existing \(out\) tail
in \([n,2048)\) is preserved. This Hoare partial-correctness theorem opens the
actual copy loop, rather than using a constructed decoder buffer.

Of \identifier{actual_rans_encode_mode2_control} and
\identifier{actual_rans_decode_mode2_control}, the encoder theorem fixes the
concrete mode-2 tables/count and canonical symbol stream while leaving the
initial state arbitrary; the decoder theorem fixes the concrete tables and
\identifier{mode2_decode_state} input. Each proves only the returned
\(bad\in\{0,1\}\). Subsequently,
\identifier{actual_rans_harness_branches_on_encoder_result} for
\cref{def:rans-actual-harness} proves the actual encoder--copy--decoder order,
decoder non-execution on failure, and exact decoder input fields in the success
branch. This conclusion contains no trace equality, decoder success, or symbol
round trip.
\end{verified}

The exact status at the end of Week~9 was as follows.

\begin{center}
\begin{tabularx}{\textwidth}{L{0.57\textwidth}L{0.25\textwidth}}
\toprule
Obligation & Status \\
\midrule
HBZ prepare/apply, leaf inverse, and bounded tail frames & \statusproved \\
Focused/full actual procedure boundary & \statusproved \\
Actual 13-symbol table certificate & \statusproved \\
Pure and reciprocal fast single-step inverse & \statusproved \\
Pure normalization-byte inverse and state bounds & \statusproved \\
Actual suffix-copy pointwise/frame theorem & \statusproved \\
Actual encoder--copy--decoder harness control & \statusproved \\
Actual encoder outer-loop trace refinement & \statuspartial \\
Actual decoder trace-consumption refinement & \statuspartial \\
Actual rANS core inverse and success witness & \statuspartial \\
Success-conditioned full-HBZ inverse & \statuspartial \\
\bottomrule
\end{tabularx}
\end{center}

The actual encoder processes symbols in reverse order, stacking normalization
bytes at a backward cursor, while the decoder consumes them from the front.
Week~9 lacked the two outer invariants establishing this backward-stack/forward-
consumption relation between actual \((i,x,off,buffer)\) and pure
\((x_j,c_j,B)\). Thus, following \cref{bdry:hbz-success}, encoder success was
not assumed from canonical HBZ alone, and
\claim{OBL-SIG-HBZ-ENCODE-DECODE} was \statuspartial. This table and
\identifier{CONTINUE-RANS} are historical verdicts at the end of Week~9.

\subsection{Step 12: actual encoder inner semantics and successful sizes}

Week~10 further decomposed the actual encoder boundary instead of moving to
the decoder. \sourcepath{Mode2RansArrayList}\allowbreak\sourcepath{Bridge.ec}
defines the array/list suffixes and segment/frame operations of
\cref{def:rans-array-list-bridge}, and proves the 1024th-suffix base equation,
one-element recurrence, canonicality, concatenation, and one-symbol
\identifier{encode_trace} recurrence.

\begin{verified}[Array/list suffixes and the actual word step]
\label{thm:rans-array-word-bridge}
\identifier{encode_trace_}\allowbreak\identifier{suffix_extension} constructs
a one-symbol trace from the array's \(i\)th symbol and \(i+1\)th suffix.
\identifier{actual_mode2_encoder_}\allowbreak\identifier{word_step_correct}
in \sourcepath{Mode2RansEncoder}\allowbreak\sourcepath{WordStep.ec} proves that,
under
\[
  0\le s<13,\quad
  1\le\operatorname{uint}(x)<\mathsf{hbz\_xmax}(s)
\]
the W32 result computed using reciprocal multiplication, shifts, and bias from
the actual mode-2 tables equals the W32 representation of pure
\identifier{hbz_fast_encode_step}. This is the word-level theorem of
\claim{OBL-RANS-ENCODER-WORD-STEP}; it is not itself a Hoare theorem for the
generated outer loop.
\end{verified}

\sourcepath{Mode2RansEncoderInnerProgress.ec} proves that pure normalization
length is 0, 1, or 2, along with division/byte addition at each step. Reusable
predicates for the actual generated loop are in
\sourcepath{Mode2RansEncoderTrace.ec}, and the direct Hoare proof is in
\sourcepath{Mode2RansEncoderActualInner.ec}.

\begin{verified}[Byte/frame steps of the actual generated inner loop]
\label{thm:rans-actual-inner}
\identifier{actual_rans_encode_inner_no_underflow} opens actual
\identifier{RansEncodeTarget.M._rans_encode} and maintains the local-step
invariant of \cref{def:rans-actual-inner}. The actual transition is
\[
  off\gets off-1;\qquad encp[off]\gets\operatorname{low8}(x);
  \qquad x\gets x\mathbin{\mathtt{>>}}8
\]
and the proof verifies preservation of the written byte segment and lower
prefix frame, and absence of underflow in the W64 decrement. The procedure's
return postcondition is
\[
  bad'\ne0\quad\lor\quad(bad'=0\land\operatorname{uint}(off')\le1020)
\]
The current phase range \(0\le k\le4\) is conservative. Identifying the actual
phases with the pure 0/1/2 normalization bytes and connecting the pre-inner-call
array to the accumulated outer trace remain \statuspartial{} under
\claim{OBL-RANS-ACTUAL-INNER-NORMALIZE}.
\end{verified}

\begin{verified}[Final serialization leaf and actual successful sizes]
\label{thm:rans-encoder-success-size}
\leavevmode\par
\sourcepath{Mode2RansEncoderSerialization.ec} proves that four actual pattern
\identifier{set8} writes equal \identifier{serialize32_le} and preserve
exterior and prefix frames. This is a serialization leaf theorem, not yet
composed into the generated procedure's full postcondition.

Separately, \identifier{actual_rans_encode_success_size_bound} retains the
failure branch of actual \identifier{_rans_encode} and proves, on success,
\[
  0\le\operatorname{uint}(off')\le1020,\qquad
  4\le1024-\operatorname{uint}(off')\le1024
\]
This is actual-procedure partial correctness for
\claim{OBL-RANS-ENCODER-SUCCESS-SIZE}. It does not prove success reachability
or termination.
\end{verified}

The exact status at the end of Week~10 was as follows.

\begin{center}
\begin{tabularx}{\textwidth}{L{0.57\textwidth}L{0.25\textwidth}}
\toprule
Obligation & Status \\
\midrule
\claim{OBL-RANS-ARRAY-LIST-BRIDGE} & \statusproved \\
\claim{OBL-RANS-ENCODER-WORD-STEP} & \statusproved\ (word level) \\
\claim{OBL-RANS-ACTUAL-INNER-NORMALIZE} & \statuspartial \\
\claim{OBL-RANS-FINAL-SERIALIZATION} & \statusproved\ (leaf theorem) \\
\claim{OBL-RANS-ENCODER-SUCCESS-SIZE} & \statusproved\ (actual partial correctness) \\
\claim{OBL-RANS-ENCODE-REFINEMENT} & \statuspartial \\
\claim{OBL-RANS-ACTUAL-SUCCESS-WITNESS} & \statuspartial \\
\claim{OBL-RANS-DECODE-REFINEMENT} and core inverse & \statuspartial \\
\claim{OBL-SIG-HBZ-ENCODE-DECODE} & \statuspartial \\
\bottomrule
\end{tabularx}
\end{center}

The remaining encoder postcondition must state simultaneously that the actual
returned suffix equals \identifier{trace_bytes} applied to
\identifier{symbol_list_of_array(symbols0)} pointwise, that its length plus the
returned offset is 1024, and that the unwritten prefix is framed. This requires
maintaining \(0\le i\le1024\), the suffix-trace state, the offset--length
equality, \identifier{segment_matches}, and \identifier{prefix_frame} together
in the outer invariant. Week~10 lacked the edge connecting one outer iteration's
actual table reads/protection/reciprocal-shift control to the already proved
inner-byte and word steps, and composing the final four stores with the
accumulated suffix.

Thus the historical Week~10 verdict was \identifier{CONTINUE-ENCODER}, and
Week~11 targeted only this single edge.

\subsection{Step 13: Week~11 actual encoder trace closure}

Week~11 directly appended the pure symbol-tail trace to Week~10's local inner
invariant and installed an exact invariant. For convenience, write
\[
\begin{aligned}
 \ell&:=\operatorname{mode2\_normalization\_len}(x_0,s),\\
 W_k&:=\operatorname{inner\_written\_suffix}(x_0,s,k),&
 B_t&:=\operatorname{encode\_trace}(tail).\mathsf{bytes}
\end{aligned}
\]
The central form is then
\[
\begin{aligned}
 I_{\mathrm{tail}}(k) \equiv{}&
 0\le k\le \ell\\
 &\land\;\operatorname{segment\_matches}\!\left(
   encp,off_0-k,W_k\mathbin{++}B_t\right)\\
 &\land\;\operatorname{prefix\_frame}(enc_0,encp,off_0-k).
\end{aligned}
\]
When the actual guard becomes false,
\(k=\operatorname{mode2\_normalization\_len}(x_0,s)\le2\) is derived. The
generated table's \(x_{\max}\), reciprocal, bias, packed complement/shift values,
and shift ladder are equated with the pure fast step through a dedicated
word bridge.

The Week~11 files have the following roles. Read each filename below with the
common prefix \sourcepath{Mode2RansEncoder}.
\begin{itemize}
  \item \sourcepath{TailInvariant.ec}: invariant initialization, one-byte
        progress, exact loop-exit characterization, and the outer one-symbol success transition;
  \item \sourcepath{GeneratedWordStep.ec}: connects generated table reads
        and word updates;
  \item \sourcepath{SerializationComposition.ec},
        \sourcepath{Finalization.ec},
        \sourcepath{GeneratedFinalization.ec}: actual final-store composition
        prepending the serialized state to the accumulated trace;
  \item \sourcepath{ActualTraceClosure.ec}: Hoare theorem directly opening
        actual inner/outer loops and final stores;
  \item \sourcepath{OuterRefinement.ec}: final corollary explicitly expanding
        the complete successful suffix.
\end{itemize}

\begin{verified}[Exact successful suffix of the actual encoder]
\label{thm:actual-encoder-trace-closure}
\identifier{actual_rans_encode_trace_closure} directly opens actual generated
\identifier{RansEncodeTarget.M._rans_encode} to prove the following
failure/success postcondition. Below, let
\[
 B_0:=\operatorname{trace\_bytes}
   (\operatorname{symbol\_list\_of\_array}(symbols_0))
\]
Then the postcondition is
\[
\begin{aligned}
 bad'\ne0\quad\lor\quad\bigl(bad'=0
 &\land 0\le off'\le1020
 \land 4\le1024-off'\le1024\\
 &\land\operatorname{segment\_matches}
   (encp',off',B_0)\\
 &\land off'+|B_0|=1024\\
 &\land\operatorname{prefix\_frame}(enc_0,encp',off')\bigr).
\end{aligned}
\]
Success is not a precondition, and the trace is computed canonically from the
actual input symbol array. \identifier{actual_rans_encode_trace_refinement}
is the final corollary expanding the same conclusion. Both theorems compiled in
the full 57-target \identifier{-no-eco} verification.
\end{verified}

The status at the end of Week~11 was as follows.

\begin{center}
\begin{tabularx}{\textwidth}{L{0.57\textwidth}L{0.25\textwidth}}
\toprule
Obligation & Status \\
\midrule
\claim{OBL-RANS-ACTUAL-INNER-NORMALIZE} & \statusproved \\
\claim{OBL-RANS-FINAL-SERIALIZATION} & \statusproved\ (actual composition) \\
\claim{OBL-RANS-ENCODE-REFINEMENT} & \statusproved\ (success-conditioned partial correctness) \\
\claim{OBL-RANS-ACTUAL-SUCCESS-WITNESS} & \statuspartial \\
\claim{OBL-RANS-DECODE-REFINEMENT} & \statuspartial \\
\claim{OBL-RANS-CORE-INVERSE} & \statuspartial \\
\claim{OBL-SIG-HBZ-ENCODE-DECODE} & \statuspartial \\
\bottomrule
\end{tabularx}
\end{center}

The end-of-Week~11 verdict was \identifier{GO-ENCODER}. This meant that the
encoder refinement edge was closed, allowing Week~12 to move to actual decoder
trace refinement. It did not mean that encoder termination, success-branch
reachability, decoder success, the HBZ round trip, or encoding zero loss had
been proved.

\subsection{Step 14: Week~12 actual decoder semantic refinement}

Week~12 did not treat existence of a pure trace as decoder success; it directly
opened both loops of actual generated
\identifier{RansDecodeTarget.M._rans_decode}. Read the filenames with the
common prefix \sourcepath{Mode2RansDecoder}.
\begin{itemize}
  \item \sourcepath{Cursor.ec}, \sourcepath{CursorSteps.ec}: reference states,
        segments, initial/one-step/final cursor boundaries, and read ranges;
  \item \sourcepath{WordStep.ec}, \sourcepath{ActualWord.ec},
        \sourcepath{GeneratedStep.ec}: connect concrete
        \identifier{symbol_words} lookups, \identifier{dsyms_words} fields,
        and actual W32/W64 updates to the mathematical decoder step;
  \item \sourcepath{Normalization.ec}: partial byte replay and exactness
        of the 0/1/2-byte guards;
  \item \sourcepath{ActualTrace.ec}: actual outer/inner loop invariants,
        initial parsing, one-symbol transitions, and exact exits;
  \item \sourcepath{TopHoare.ec}: top-level Hoare theorem directly targeting
        the actual generated procedure.
\end{itemize}

The reference boundaries for the expected symbol list \(s\) are
\[
\begin{gathered}
 c_0=4,\qquad c_{i+1}=c_i+|\beta_i|,\qquad
 c_{1024}=|\mathsf{trace\_bytes}(s)|,\\
 x_0=\operatorname{ParseLE32}(\mathsf{trace\_bytes}(s)[0..4)),\qquad
 x_{1024}=2^{23}.
\end{gathered}
\]
\identifier{generated_decoder_parse32_from_trace} connects the actual 4-byte
parsing in the first equation, and generated-table theorems match the symbol
selected by \(x_i\bmod1024\) and the reduced state to the mathematical step.
The inner invariant maintains
\[
 k=\operatorname{uint}(off)-c_i,\qquad
 x=\mathsf{decoder\_replay\_prefix}(x_{i+1},s_i,k)
\]
The generated guard equals \(k<|\beta_i|\), and \(|\beta_i|\le2\), so the
actual loop consumes exactly 0, 1, or 2 bytes.

\begin{verified}[Exact-trace partial correctness of the actual decoder]
\label{thm:actual-decoder-trace-closure}
The preconditions of \identifier{actual_rans_decode_trace_refinement} are a
canonical 1024-symbol array, exact \identifier{trace_bytes} computed from that
array, size \(4\le n\le1024\), state fields \((1024,n,13)\), the two actual
mode-2 decoder tables, and a pointwise trace-read relation. Neither success nor
output equality is a precondition. The postcondition of every terminating
execution is
\[
\begin{aligned}
  bad'&=0,& off'&=n,\\
  decoded'[0..1024)&=expected[0..1024),&
  decoded'[1024..2048)&=decoded_0[1024..2048)
\end{aligned}
\]
The internal exit fact \(x=2^{23}\) discharges the actual final rejection check,
but is not exposed in the return ABI. This theorem directly targets the actual
generated procedure; it is not a wrapper or pure-decoder theorem.
\end{verified}

The status at the end of Week~12 was as follows.

\begin{center}
\begin{tabularx}{\textwidth}{L{0.57\textwidth}L{0.25\textwidth}}
\toprule
Obligation & Status \\
\midrule
Decoder cursor/table/word-step bridges & \statusproved \\
\claim{OBL-RANS-DEC-NORMALIZE} & \statusproved\ (actual-loop connection) \\
\claim{OBL-RANS-DECODE-REFINEMENT} & \statusproved\ (exact-trace partial correctness) \\
\claim{OBL-RANS-CORE-INVERSE} & \statuspartial \\
\claim{OBL-RANS-ACTUAL-SUCCESS-WITNESS} & \statuspartial \\
\claim{OBL-SIG-HBZ-ENCODE-DECODE} & \statuspartial \\
\bottomrule
\end{tabularx}
\end{center}

Full 65-target \identifier{-no-eco} verification passed, and the Week~12 verdict
was \identifier{GO-DECODER}. The single Week~13 edge was composition transporting
the encoder's \identifier{segment_matches} and copy-helper
\identifier{slice_eq} to the decoder's pointwise exact-trace input inside actual
\identifier{Mode2RansActualHarness.run}. This composition, actual encoder-success
reachability, termination of the two loops, full-HBZ wrappers, and encoding
zero loss do not follow from the Week~12 theorem.

\subsection{Step 15: Week~13 actual rANS core inverse composition}

Week~13 closed Week~12's last unconnected edge with two authored theories.
\sourcepath{Mode2RansCore}\allowbreak\sourcepath{CompositionBridge.ec} supplies
the following postcondition adapters.

\begin{itemize}
  \item \identifier{copied_suffix_is_exact_trace}: derives the exact trace at
        offset 0 of the copied buffer from actual encoder
        \identifier{segment_matches} and actual copy \identifier{slice_eq}.
  \item \identifier{segment_matches_implies_exact_}\allowbreak
        \identifier{decoder_segment_input}: expands the global trace segment
        into pointwise decoder reads in each normalization segment.
  \item \identifier{actual_decoder_input_from_}\allowbreak
        \identifier{configured_trace}: constructs Week~12's full decoder-input
        predicate from the three actual state stores and the copied trace.
  \item \identifier{encoder_success_size_word_bridge}: proves non-wrapping
        \identifier{W64} subtraction and equality with the integer trace size
        from the actual successful-offset range.
\end{itemize}

\begin{verified}[Success-conditioned core inverse of three actual procedures]
\label{thm:actual-rans-core-inverse}
\identifier{actual_rans_encode_copy_decode_inverse} in
\sourcepath{Mode2RansCoreActualInverse.ec} directly targets
\identifier{Mode2RansActualHarness.run}. Its only preconditions are bindings
of six initial arguments and canonical \identifier{mode2_hbz_symbol_stream};
it does not require encoder success. Every terminating execution satisfies one
of the following two branches.
\[
\begin{aligned}
  bad_{\mathrm{enc}}\ne0&\;\Longrightarrow\;
    \neg\mathsf{decoder\_ran}\land
    decoded'=decoded_0\land state'=state_0,\\
  bad_{\mathrm{enc}}=0&\;\Longrightarrow\;
    \mathsf{decoder\_ran}\land bad_{\mathrm{dec}}=0
    \land off_{\mathrm{dec}}=\mathsf{encoded\_size}\\
  &\hspace{7.2em}\land\;
    decoded'[0..1024)=symbols[0..1024).
\end{aligned}
\]
The success branch also preserves decoder input fields
\((1024,\mathsf{encoded\_size},13)\), the output \([1024,2048)\) tail frame,
the actual encoder's exact suffix, and the initial prefix frame. Thus there is
no circular assumption repeating the copied trace or decoder results as premises.
\end{verified}

The theorem's premises are satisfiable by a canonical all-zero symbol array,
but its conclusion branches on actual encoder failure/success, so success-branch
reachability does not follow. The composition module is also a proof-authored
harness directly calling three fixed generated procedures, not the production
\identifier{_encode_hb_z1_full}/\identifier{_decode_hb_z1_full} wrappers themselves.

The completion status for Week~13 was as follows.

\begin{center}
\begin{tabularx}{\textwidth}{L{0.57\textwidth}L{0.25\textwidth}}
\toprule
Obligation & Status \\
\midrule
Copy-trace/pointwise decoder-input/W64 size adapters & \statusproved \\
\claim{OBL-RANS-CORE-INVERSE} & \statusproved\ (success-conditioned partial correctness) \\
\claim{OBL-RANS-ACTUAL-SUCCESS-WITNESS} & \statuspartial \\
Encoder/decoder termination & \statuspartial \\
\claim{OBL-SIG-HBZ-ENCODE-DECODE} & \statuspartial\ (actual full-wrapper composition) \\
\bottomrule
\end{tabularx}
\end{center}

Full 67-target \identifier{-no-eco} verification and independent theorem-surface
checks passed, and the Week~13 verdict was \identifier{GO-CORE}. The single
Week~14 implementation edge was success-conditioned composition of the already
closed prepare/apply inverse and this core inverse at the actual full-HBZ
wrapper boundary. Actual success witnesses, termination, the h codec, the full
1474-byte codec, and encoding zero loss do not follow from the Week~13 theorem.

\subsection{Step 16: Week~14 production full-HBZ wrapper composition}

Week~14 closed the production boundary left by Week~13 with five authored
theories. \sourcepath{Mode2HbzInternal}\allowbreak\sourcepath{Boundaries.ec}
connects the prepare, rANS encode/decode, and apply procedures in full-wrapper
and existing focused extractions by exact equivalence. The following two
adapters connect the coefficient and symbol boundaries.

\begin{itemize}
  \item \identifier{prepared_hbz_implies_mode2_symbol_stream}: transports
        actual prepare results on canonical HBZ to the Week~13 core inverse's
        13-symbol-stream premise.
  \item \identifier{decoded_prefix_preserves_prepared_hbz}: returns the symbol
        prefix recovered by the actual decoder to the apply theorem's
        prepared-HBZ premise.
\end{itemize}

\sourcepath{Mode2HbzFullEncodeTrace.ec} and
\sourcepath{Mode2HbzFullDecodeInverse.ec} directly target the actual full
encoder and decoder, respectively. The former concludes actual prepare-local
symbols and the exact encoded trace, size, and suffix frame; the latter
concludes \(bad=0\), the original 1024 coefficients, and the coefficient tail
frame from that trace. Both results are composed in the direct harness of
\sourcepath{Mode2HbzFullActualInverse.ec}.

\begin{verified}[Production full-HBZ wrapper inverse]
\label{thm:production-hbz-wrapper-inverse}
The boundary lift in \sourcepath{Mode2HbzSignatureBoundaryLift.ec} proves exact
equality of result tuples and the \identifier{decoder_ran} flag between the
following two executions.
\begin{itemize}[leftmargin=2em]
  \item production: \identifier{SignaturePackUnpackHbzFullHarness.run}
  \item focused: \identifier{HbzFullActualHarness.run}
\end{itemize}
The exact-equivalence theorem is
\identifier{signature_pack_unpack_hbz_full_actual_exact}. Its Hoare corollary
\identifier{signature_pack_unpack_hbz_full_inverse_mode2} assumes only canonical
HBZ and initial argument bindings and shows the following.
\[
\begin{array}{ll}
size=0 &: \neg\mathsf{decoder\_ran}\text{ and initial-buffer preservation},\\
size\ne0 &: 4\le\operatorname{uint}(size)\le1024,
\ \mathsf{decoder\_ran},\ bad=0,\\
& decoded[0..1024)=hbz_0[0..1024)\text{ and two tail frames}.
\end{array}
\]
The success branch also supplies actual prepare's
\identifier{prepared_symbols} and exact \identifier{mode2_trace_bytes}
witnesses. This is not circular composition: neither success nor decoder output
is required as a precondition.
\end{verified}

The completion status for Week~14 was as follows.

\begin{center}
\begin{tabularx}{\textwidth}{L{0.57\textwidth}L{0.25\textwidth}}
\toprule
Obligation & Status \\
\midrule
Full/focused internal-procedure equivalence and coefficient--symbol adapters & \statusproved \\
Direct refinement of actual full-HBZ encoder/decoder & \statusproved \\
\claim{OBL-SIG-HBZ-ENCODE-DECODE} & \statusproved\ (success-conditioned partial correctness) \\
\claim{OBL-RANS-ACTUAL-SUCCESS-WITNESS} & \statuspartial \\
Encoder/decoder termination & \statuspartial \\
\claim{OBL-SIG-H-ENCODE-DECODE} & \statusspecified \\
\bottomrule
\end{tabularx}
\end{center}

Full 72-target \identifier{-no-eco} verification, including the five new
full-HBZ theories, passed, and the Week~14 verdict was \identifier{GO-HBZ}.
The next single implementation obligation was to show within EasyCrypt that
all-zero HBZ coefficients---an all-6 symbol stream after prepare---reach the
actual nonzero-size success branch. Before this witness, the result is not
extended to unconditional round trips, the (h) codec, the full 1474-byte codec,
or encoding zero loss.

\subsection{Step 17: Week~15 actual success postcondition for fixed all-6 input}

Week~15 used two authored theories and a minimal strengthening of the existing
encoder-closure theorem.
\sourcepath{Mode2RansAllSix}\allowbreak\sourcepath{Budget.ec} supplies the
following connections.

\begin{itemize}
  \item Constructs canonical HBZ and the actual prepare's all-6 symbol prefix
        from actual 32-bit reads of bytewise \identifier{zero_hbz}.
  \item Calculates that symbol 6 emits at most one normalization byte for state
        \(x<2^{31}\), and that the first four steps emit 0 bytes.
  \item Hence the 1024-symbol normalization trace has at most 1020 bytes, and
        the full trace including the final 4 state bytes fits in the fixed
        1024-byte workspace.
\end{itemize}

The existing \sourcepath{Mode2RansEncoderActual}\allowbreak
\sourcepath{TraceClosure.ec} was strengthened.
\identifier{actual_rans_encode_}\allowbreak\identifier{failure_trace_cause}
transports from the existing loop invariant the fact that actual failure
requires a normalization trace longer than 1020 bytes. Contradiction with the
pure upper bound closes the following actual theorem.

\begin{verified}[Fixed all-6 production HBZ success]
\label{thm:all-six-production-success}
\identifier{actual_rans_encode_all_six_success} in
\sourcepath{Mode2RansActual}\allowbreak\sourcepath{SuccessWitness.ec} directly
targets actual \identifier{RansEncodeTarget.M._rans_encode}, without assuming
success or buffer fit as preconditions. On terminating executions, it gives
\[
 bad=0,\qquad 0\le off\le1020,\qquad
 4\le1024-off\le1024
\]
and the exact all-6 trace and prefix frame.
\identifier{actual_encode_hb_z1_full_zero_success} and
\identifier{signature_pack_hbz_zero_success_mode2} lift this conclusion to the
actual full encoder and production packer. Finally,
\identifier{signature_pack_unpack_hbz_zero_success_mode2} concludes nonzero
size, decoder execution, \(bad=0\), recovery of 1024 zero coefficients, and two
tail frames in the production harness on fixed \identifier{zero_hbz} input.
\end{verified}

The completion status for Week~15 was as follows.

\begin{center}
\begin{tabularx}{\textwidth}{L{0.57\textwidth}L{0.25\textwidth}}
\toprule
Obligation & Status \\
\midrule
Zero-HBZ prepare-to-all-6 and 1020-byte budget & \statusproved \\
\claim{OBL-RANS-ACTUAL-SUCCESS-WITNESS} &
  \statusproved\ (fixed-input Hoare partial correctness) \\
Fixed zero-HBZ production round trip & \statusproved\ (Hoare partial correctness) \\
Fixed-input and general rANS termination/losslessness & \statuspartial \\
\bottomrule
\end{tabularx}
\end{center}

Full 74-target verification, including the two new theories, passed, and the
Week~15 verdict was \identifier{GO-WITNESS}. This is not
\identifier{GO-WITNESS-LOSSLESS}. Execution existence, termination,
probability-1 success, and success on every canonical HBZ input do not follow.
The subsequent one-week minimum paper scope superseded the original (h)-codec
recommendation, deferring (h).

\subsection{Step 18: Week~16 MINCORE KeyGen snapshot}

The first Week~16 task was transparent composition of two actual helpers,
rather than the full KeyGen or API. \sourcepath{Mode2KeygenSnapshotAlgebra.ec}
proves, as pure integer theorems, for canonical residue \(r\),
\[
  r=2\operatorname{high}(r)+\operatorname{low}(r),\qquad
  -1\le\operatorname{low}(r)\le1
\]
and mod-\(2q\) congruence of raw sums/residues.
\identifier{ActualM23MatrixFinalizeSnapshot.run} in
\sourcepath{Mode2KeygenCoreEquation.ec} directly calls actual
\identifier{_kp_m23_matrix} and \identifier{_keypair_finalize_m23} in sequence,
preserving the intermediate \(bp\) as \identifier{pre_bp}.

\begin{verified}[Actual matrix/finalizer semantic snapshot]
\label{thm:actual-keygen-snapshot}
\identifier{actual_m23_matrix_finalize_semantic_snapshot} assumes only
input/matrix bounds and initial array bindings, and jointly concludes the following.
\begin{itemize}
  \item Actual matrix-helper \identifier{mode2_output_repr_bound16} and
        \identifier{mode2_ntt_repr_bound24};
  \item Exact array/HAETAE semantics of the actual finalizer and
        \(r=2b_1+b_0,\ -1\le b_0\le1,\ s'_2=s_2-b_0\);
  \item For every active coefficient,
        \(2(a-2b_1)+2\,pre\_bp+2s'_2\equiv0\pmod{2q}\);
  \item Initial-to-snapshot-to-final, NTT scratch, and \(s_2\) tail frames.
\end{itemize}
Neither final outputs nor the desired KG predicate appear as preconditions.
\end{verified}

However, \identifier{pre_bp} is only a full-NTT/pointwise-product/inverse-NTT
snapshot represented by \identifier{KeygenM23ArithmeticSpec.output_row}.
There was no theorem identifying it with the paper's polynomial product
\(A_{\mathrm{gen}}s_{\mathrm{gen}}\), so the status at that time was as follows.

\begin{center}
\begin{tabularx}{\textwidth}{L{0.57\textwidth}L{0.25\textwidth}}
\toprule
Obligation & Status \\
\midrule
Actual matrix/finalizer snapshot semantics & \statusproved \\
KG-2-style actual coefficient low/high decomposition & \statusproved \\
Adjusted \(s_2\) component of KG-4 / full KG-4 & \statusproved\ /\ \statuspartial \\
Faithful KG-1, KG-3, paper \(As=qj\bmod2q\) & \statusblocked \\
\claim{OBL-MINCORE-SIGN}, \claim{OBL-MINCORE-VERIFY},
\claim{OBL-MINCORE-COMPOSITION} & \statuspartial\ /\ \statuspartial\ /\ \statusspecified \\
\claim{OBL-SIG-H-ENCODE-DECODE} & \statusdeferred \\
\bottomrule
\end{tabularx}
\end{center}

The pre-change 76-target baseline, including the two KeyGen snapshot theories,
passed full verification. The first task's historical verdict was
\identifier{CONTINUE-KG} and \identifier{BLOCKED-KG-NTT-MUL}; the continuation
audit below updated it to the final stop verdict.

\subsection{Step 19: Week~16 KG-NTT-MUL audit and stop verdict}

The follow-up task made the last representation edge of the existing actual
NTT chain explicit in the source, without assuming the desired security-model
equality. \identifier{output_row_from_mode2_ntt_words} in the 77th manifested
target, \sourcepath{Mode2KeygenNttMulBridge.ec}, states, under the transformed
secret-key representation,
\[
 \mathsf{output\_row}
 =\operatorname{array256\_mont}
   (\operatorname{full\_invntt}(\operatorname{pointwise\_row\_words}(\cdots)))
\]
and \identifier{output_row_repr_from_mode2_ntt_words} transports that equality
to the existing \identifier{wide_slice_repr_bound} contract.
\identifier{actual_m23_matrix_snapshot_rows_explicit} directly applies this
representation to the postconditions of the two active rows, as a consequence
of the existing actual two-call snapshot theorem. None of the three theorems
mentions or assumes \identifier{Rq.\&*},
\identifier{HAETAE_Algebra.poly_mul},
\(A_{\mathrm{gen}}s_{\mathrm{gen}}\), or KG predicates. The file, including
these three theorems and the direct harness import, passed individual fresh
\identifier{-no-eco} compilation.

\begin{boundary}[Two missing KG-NTT-MUL leaves at paper freeze]
First, there was no odd-root orthogonality/full-NTT convolution theorem showing,
with the actual Montgomery factor,
\[
 \operatorname{mont}
  (\operatorname{full\_invntt}(\widehat a\odot\operatorname{full\_ntt}(p)))
 =\operatorname{full\_invntt}(\widehat a)\mathbin{\star_{\!Rq}}p
\]
Second, there was no adapter transporting products of native \identifier{Rq.poly}
arrays to the security model's integer-list
\identifier{poly_mul}/\identifier{matrix_vec_mul}. Legacy \identifier{Rq.ntt}
differs from the actual full transform and cannot fill this gap. Thus renaming
the snapshot's \identifier{pre_bp} as
\(A_{\mathrm{gen}}s_{\mathrm{gen}}\) is circular.
\end{boundary}

The frozen verdict is \identifier{STOP-KG-NTT}. KG-2/finalization is
\statusproved; only the adjusted \(s_2\) component of KG-4 is \statusproved;
faithful KG-1/KG-3, full KG-4, and the paper equation were not proved. The KeyGen
frontier is frozen at this boundary. The 78th target,
\identifier{Mode2SignAcceptedCore}, subsequently compiled direct calls to the
three actual helpers and accepted-branch control. However, the full
challenge-semantics leaf is absent, so the verdict is
\identifier{STOP-SIGN-CHAL-MODE2}, and S-1--S-7 are unproved. The four subsequent
Verify targets compiled sequential control of the five actual helpers, exact
machine-word expressions for V-1/V-2/V-5/V-6, the W64 norm gate, and tail/mismatch
word expressions. With no matrix CRT semantic bridge, the verdict is
\identifier{STOP-VERIFY-MATRIX-CRT}; at freeze time, the first leaf for resuming
work was \identifier{rq_mul_coeff_foldr_to_bigi}, followed by the 2-by-4 NTT/CRT
procedure leaves.

\subsection{Step 20: paper freeze and post-freeze KeyGen semantic alignment}
\label{subsec:post-freeze-keygen}

Commit \identifier{4bde1207e2a290952d491d5b41c4bd773081e3c4} added no new
theorems; it aligned the 82 authored targets, claim ledger, theorem graph,
reproduction procedures, and paper PDF at the same \identifier{PAPER-FROZEN}
boundary. The frozen PDF's SHA-256 is
\hashsplit{be935948028829556951863b44ceb2e6}{c5b2037820991b3a43a4b461b036e53d}.
Subsequent work proceeded in the external NTT tree, the existing 77th KeyGen
target, and \sourcepath{post-freeze/} auxiliary targets without retroactively
changing this paper claim matrix.

The first auxiliary layer closes the pure NTT and representation edges blocked
at paper freeze.
\sourcepath{haetae-ntt-verify/easycrypt/}\allowbreak
\sourcepath{NTTFullSpectralAction.ec} proves
\identifier{rq_mul_coeff_foldr_to_bigi},
\identifier{full_ntt_montgomery_spectral_action}, and the general
\identifier{full_ntt_montgomery_row_product}. It does not use legacy
\identifier{Rq.ntt}, which differs from the actual full transform.
\sourcepath{NTTMode2Row}\allowbreak\sourcepath{Specializations.ec} specializes
these results to the mathematical 2-by-3 KeyGen and 2-by-4 Verify row products.

\sourcepath{RqHAETAEBridge.ec} introduces the coefficientwise representation
relation \identifier{rq_poly_repr}. Its only premises are well-formedness of
HAETAE polynomials and equality of active coefficients; they do not include the
desired convolution or KG equation. The following theorems were freshly compiled.

\begin{itemize}
  \item \identifier{rq_poly_add_repr}: representation preservation for
        \identifier{Rq.\&+} and \identifier{HAETAE_Algebra.poly_add};
  \item \identifier{rq_poly_mul_repr}: representation preservation for checked
        negacyclic \identifier{Rq.\&*} and \identifier{HAETAE_Algebra.poly_mul};
  \item \identifier{rq_poly_dot3_repr}: transports the sum of three products
        to \identifier{HAETAE_Algebra.poly_dot};
  \item \identifier{keygen_mode2_row_product} and
        \identifier{mode2_row_product_haetae_dot}: connect actual KeyGen NTT
        output to the native Rq row product and HAETAE three-term dot product.
\end{itemize}

Thus the two pure odd-root/full-NTT and Rq--HAETAE leaves missing at paper freeze
are no longer globally unresolved obligations. However, this alone does not
automatically identify actual matrix words, security-side \(A\), secret/error
packaging, and the paper's mod-\(2q\) equation.

The second auxiliary layer aligns the deterministic semantics of terminating
actual mode-2 \identifier{avec} expansion with the paper objects.
\sourcepath{post-freeze/KgActualAvecQjSemantics.ec} proves the following.

\begin{itemize}
  \item The two row nonces are 515 and 516, respectively, and the SHAKE128
        input appends two nonce bytes to the 32-byte seed;
  \item Domain byte 31, final padding byte 128, 168-byte squeeze blocks, and
        the 672-byte size of the first four blocks match actual framing;
  \item Adjacent bytes are read as little-endian 16-bit candidates, only
        values below \(64513\) are accepted, and accepted coefficients are
        stored without further modular reduction;
  \item The two 256-coefficient rows obtained through this process are the
        paper-level mode-2 vector \(a\).
\end{itemize}

In the paper's \(j=(1,0)^T\), the first \(1\) is the identity constant
polynomial, so \(qj\) adds \(q=64513\) only to coefficient~0 of row~0, with all
other coefficients zero. \identifier{paper_qj_row0_coeff0},
\identifier{paper_qj_row0_other}, and \identifier{paper_qj_remaining_row}
separate these cases.

The two auxiliary theorems in
\sourcepath{post-freeze/}\allowbreak
\sourcepath{KgActualAvecQjComposition.ec},
\begin{center}
\small
\identifier{mode2_sampler_facts_actual_paper_a},\qquad
\identifier{actual_snapshot_zero_adds_paper_qj}
\end{center}
compose these two paper objects with actual sampler facts and the existing
snapshot mod-\(2q\) equation. The top-level Hoare theorem
\begin{center}
\small
\identifier{checked_mode2_parent_m23_finalize_actual_avec_qj}
\end{center}
exposes both results for terminating executions of the checked sampler, actual
matrix helper, and finalizer. Sampler termination and distribution are not
conclusions.

\begin{boundary}[\identifier{STOP-KG-AVEC-QJ}]
The existing security model's \identifier{haetae_mode23_qj_vector} is neither
actual ExpandVecA nor paper \(qj\), but a synthetic arithmetic generator using
tag 401 in \identifier{public_key_seed_polyveck}. Compiled theorems calculate
coefficient~0 of its rows~0/1 as 401 and 403, respectively, for every seed,
whereas paper \(qj[0][0]=64513\).
\identifier{security_named_qj_is_not_paper_qj} directly proves this inequality.
An independent FIPS-202 all-zero-seed trace also gives 44985 for the actual first
coefficient of nonce-515 row~0, differing from the synthetic value 401.
Assuming the desired identity would concretely require \(44985=401\) and is
therefore unsound.
\end{boundary}

Thus the latest minimal repair obligation is
\identifier{OBL-KG-SECURITY-EXPANDVECA-SEMANTICS}. The security-side
\identifier{public_rounding_vector_seed} must be replaced or refined with actual
SHAKE128, little-endian parsing, \(<64513\) rejection, and nonce 515/516
semantics; the paper \(qj\) correction must remain a separate mod-\(2q\)
object, and an actual-array representation must then be proved. Even afterward,
the current HAETAE algebra's mode-2 width-4/mod-\(q\) representation must be
connected to the paper augmented matrix's width-\(1+3+2=6\)/mod-\(2q\)
packaging. This step therefore claims neither KG-1--KG-4,
\(As=qj\pmod{2q}\), sampler distribution, termination, nor
\(\delta_{\mathrm{KeyGen}}=0\).
